\documentclass[11pt,a4paper]{article}
% Estilo comum dos PDFs individuais da Sprint 1.
\usepackage[utf8]{inputenc}
\usepackage[T1]{fontenc}
\usepackage[brazilian]{babel}
\usepackage[a4paper,margin=2.35cm]{geometry}
\usepackage{xcolor}
\usepackage{hyperref}
\usepackage{enumitem}
\usepackage{listings}
\usepackage{parskip}
\usepackage{needspace}
\definecolor{azul}{RGB}{20,63,112}
\hypersetup{colorlinks=true,linkcolor=azul,urlcolor=azul}
\setlist[itemize]{leftmargin=*,itemsep=3pt}
\lstset{basicstyle=\ttfamily\small,frame=single,breaklines=true,columns=fullflexible}
\sloppy
\newcommand{\codigo}[1]{\texttt{#1}}
\newcommand{\titulo}[2]{\begin{center}{\LARGE\bfseries #1\par}\vspace{4pt}{\large #2\par}\vspace{8pt}\rule{0.8\textwidth}{0.4pt}\par\vspace{5pt}{\small Matheus Moreli Durão Silva -- Matrícula 12621CBS217\par Cibersegurança - Bacharelado - Noturno\par}\end{center}}
\newcommand{\cardfonte}[1]{\par\smallskip\noindent\textit{Referência: }\href{#1}{Trello - enunciado e checklist do card}.}

\begin{document}\titulo{S1\_04 - Primeiro programa em Python}{Saída no terminal com \codigo{print}}
\section*{O que o card pede}
O entregável é um programa ``Hello, World!'' executado corretamente. O objetivo não é a complexidade: é confirmar que Python está instalado, que o arquivo \codigo{.py} é reconhecido e que o aluno consegue executar um script pelo terminal.
\section*{Como \codigo{print} funciona}
\codigo{print()} envia uma representação de um valor para a tela. Texto é escrito entre aspas: \codigo{print("Olá")}. O argumento \codigo{sep} define o separador entre valores e \codigo{end} muda o final da linha. A quebra de linha comum é \codigo{\textbackslash n}: \codigo{print("linha 1\textbackslash nlinha 2")}. Quando uma variável precisa aparecer dentro de um texto, uma f-string torna a leitura simples: \codigo{print(f"Ativo: \{nome\}")}. As chaves são substituídas pelo valor da variável.
\begin{lstlisting}
nome = "Servidor de arquivos"
print("Hello, World!")
print(f"Exemplo de ativo: {nome}")
\end{lstlisting}
\section*{Aplicação no projeto}
Todo diálogo do \codigo{app.py} usa \codigo{print}: título do menu, listas de ativos, confirmações e erros. A diferença é que, no projeto final, o texto mostrado depende dos dados guardados. O primeiro programa ensina o mecanismo básico de comunicação do programa com quem o utiliza.
\section*{Como executar e conferir}
No terminal, entrar na pasta do arquivo e executar \codigo{python S1\_04.py}. A saída deve conter ``Hello, World!'', o nome e o curso. Se nada aparecer, conferir o nome do arquivo, o comando usado e a instalação do interpretador.
\section*{Checklist demonstrado}
Instalação/configuração do Python, estrutura básica de um script, comando \codigo{print}, execução e primeiro programa foram revisados e praticados.
\cardfonte{https://trello.com/c/YeugQuYt}
\end{document}
